\documentclass[11pt]{article}

\usepackage[margin=0.88in]{geometry}
\usepackage{amsmath}
\usepackage{amssymb}
\usepackage{array}
\usepackage{booktabs}
\usepackage{enumitem}
\usepackage{hyperref}
\usepackage{listings}
\usepackage{longtable}
\usepackage{microtype}
\usepackage{xcolor}
\usepackage{url}

\hypersetup{
  colorlinks=true,
  linkcolor=blue!55!black,
  citecolor=blue!55!black,
  urlcolor=blue!55!black
}

\lstset{
  language=Python,
  basicstyle=\ttfamily\small,
  columns=fullflexible,
  frame=single,
  breaklines=true,
  showstringspaces=false,
  keywordstyle=\bfseries\color{blue!60!black},
  commentstyle=\itshape\color{green!45!black}
}

\newcommand{\tg}{\textsc{TensorGuard}}
\newcommand{\safe}{\textsc{Safe}}
\newcommand{\unknown}{\textsc{Unknown}}
\newcommand{\unsafe}{\textsc{Unsafe}}
\newcommand{\rp}{\textsc{Refuted-Proof}}
\newcommand{\todoartifact}[1]{\texttt{#1}}

\title{\tg: An Auditable Static Safety Net for PyTorch Tensor Programs}
\author{halley-labs}
\date{Generated from repository artifacts, June 2026}

\begin{document}
\maketitle

\begin{abstract}
\tg{} is a zero-annotation static verifier for PyTorch \texttt{nn.Module}
architectures and adjacent training/deployment code.  It checks tensor shape,
dtype, device, phase, gradient-flow, stride/layout, probabilistic batch/event
contracts, complex view/FFT shape-dtype rules, named-axis alignment,
\texttt{vmap} batched-dimension transfers, sparse layout/blocksize invariants,
geometric sampler grids, and library-specific contracts before
a forward pass runs, then reports a
three-valued verdict:
\safe{}, \unknown{}, or \unsafe{}.  The implementation is intentionally more
than a prototype: it ships a formal soundness contract, a verifiable-fragment
specification, operator confidence tags, SARIF/JSON/JUnit/GitHub reporters,
pre-commit and pytest integrations, a safe AST/fx-only loading path for
untrusted model files, a community stub governance process, and an upstream
PyTorch hook proposal.

The evidence base is also first-class.  A frozen real benchmark corpus, a
2,704-item public GitHub bug snapshot, a 227-case generated corpus, a blind
held-out split, fuzzers, mutation tests, live PyTorch dispatcher comparisons,
statistical tests, numeric-claim audits, and a one-command reproducibility
capsule all live in the repository.  On the frozen 16-case real benchmark
corpus, \tg{} reports TP=8, FP=0, TN=8, FN=0; on 80 clean sound-mode models and
200 random clean fuzzed modules it reports zero false positives; on 281 genuine
injected faults it reports perfect recall; and on 2,000 dispatcher-generated
modules it agrees exactly with eager PyTorch with no abstentions.  The formal
side includes a Lean 4 library whose 371 public theorems are discovered and
axiom-audited by tests to depend only on the trusted Lean kernel axioms and not
on \texttt{sorryAx}, plus parity checks against the Python verifier, Z3, cvc5,
and the live torch dispatcher.
\end{abstract}

\tableofcontents
\newpage

\section{Why this tool exists}

Tensor bugs are unusually expensive because a model can be syntactically valid,
type-check in Python, instantiate successfully, and still crash only after a
particular batch shape, device placement, train/eval mode, attention mask, or
export path reaches the wrong operator.  Runtime tests catch some of these
bugs, but only when the test constructs the right input and executes the failing
path.  Smoke tests miss latent eval-only branches, CUDA-only device mismatches
on CPU machines, silent gradient breaks, and shape constraints hidden inside
library calls such as \texttt{einops.rearrange},
\texttt{scaled\_dot\_product\_attention}, \texttt{grid\_sample}, or PyTorch
named-tensor \texttt{align\_to}.

\tg{} attacks that failure mode with static verification.  Given model source
or a live module, it extracts a symbolic tensor program, propagates abstract
tensor facts through operator transfer functions, emits solver obligations, and
surfaces either a proof of safety within the declared fragment, a concrete
refutation, or an explicit abstention.  The design choice that pervades the
repository is honesty: when the tool cannot soundly prove a program safe, the
strict mode returns \unknown{} rather than a silent pass; when a benchmark number
depends on CUDA, HuggingFace downloads, or a Lean toolchain, the numeric audit
labels it environment-qualified instead of pretending it was regenerated in the
ordinary CPU-only path.

\paragraph{What is new in this repository.}
\tg{} combines a production-facing developer tool with a paper-grade evidence
machine.  The same tree contains the verifier, its soundness contract, its
generated specs, its benchmark corpora, the comparison baselines, the
statistical tests, the dashboard, the artifact-evaluation capsule, and the
formal Lean audit.  This paper summarizes those assets as a coherent system
rather than as a sequence of roadmap checkboxes.

\paragraph{Core claims.}
\begin{enumerate}[leftmargin=2em]
  \item \textbf{No-annotation PyTorch verification.}  The common path is
  \texttt{tensorguard verify model.py}; the verifier infers shapes and other
  tensor facts from constructors, \texttt{forward}, dataflow, and optional input
  specs.
  \item \textbf{Sound abstention boundary.}  Strict \texttt{sound} mode returns
  \safe{} only inside the verifiable fragment and returns \unknown{} for
  unsupported constructs.  More permissive modes trade abstention for legacy
  recall while preserving the reported bug set.
  \item \textbf{Cross-domain bug coverage.}  Shape, device, dtype, phase, and
  gradient-flow checks are measured separately.  Device and gradient each catch
  real bugs missed by a shape-only view; phase is reported as diagnostic-only
  where it does not refute.
  \item \textbf{Repository-anchored evaluation.}  Headline numbers are generated
  from committed artifacts and checked by \path{reproducibility/audit_numeric_claims.py}.
  The audit currently reports 12 verified claims, one regime-qualified claim,
  and one environment-qualified claim.
  \item \textbf{Machine-checked core obligations.}  Lean modules model the
  central operator rules, reduced products, CEGAR bounds, fragment modes, and
  SMT encodings; tests discover and axiom-audit all 371 public theorems.
\end{enumerate}

\section{User-visible capabilities}

The repository is not just a verifier kernel.  It exposes an adoption surface
that lets users run \tg{} locally, in CI, in editors, and in upstream-style
hooks.

\subsection{Command-line and Python API}

The CLI entrypoint is \texttt{tensorguard}.  Its central command verifies a file
or module and supports independent check toggles, soundness modes, and machine
readable output.  The public Python API mirrors the CLI through
\texttt{verify\_architecture} and \texttt{verify\_module}.  Specialized helpers
cover einops, SDPA, distributions, named tensors, complex FFT/view contracts,
\texttt{vmap} batched-dimension transfer, sparse tensor layouts, and
\texttt{grid\_sample}/\texttt{affine\_grid}.  A typical
workflow is:

\begin{lstlisting}
import torch.nn as nn
from src.api import verify_architecture

class Bad(nn.Module):
    def __init__(self):
        super().__init__()
        self.fc1 = nn.Linear(8, 16)
        self.fc2 = nn.Linear(99, 4)

    def forward(self, x):
        return self.fc2(self.fc1(x))

result = verify_architecture(Bad, input_shapes={"x": (2, 8)},
                             soundness_mode="sound")
assert result.verdict == "UNSAFE"
\end{lstlisting}

The CLI can emit human-readable diagnostics as well as JSON, SARIF 2.1.0,
JUnit-XML, and GitHub annotation formats.  SARIF is validated for GitHub Code
Scanning.  The pytest plugin lets projects fail tests on unsafe modules.  The
pre-commit hook and GitHub Action paths allow incremental rollout, and baseline
or suppression support lets teams adopt the tool without blocking on a legacy
backlog.

\subsection{Soundness modes and exits}

\tg{} exposes three modes:

\begin{center}
\begin{tabular}{p{0.16\linewidth} p{0.72\linewidth}}
\toprule
\textbf{Mode} & \textbf{Contract} \\
\midrule
\texttt{sound} & Report \safe{} only when the analyzed region is fully inside
the verifiable fragment, uses no heuristic-tagged operators, and has no opaque
layers.  Unsupported code returns \unknown{}; the CLI exits 2 on \unknown{} so
CI can gate on abstentions. \\
\texttt{balanced} & Default mode.  Preserves legacy usability by abstaining on
opaque layers but otherwise returning \safe{} when no bug is found. \\
\texttt{heuristic} & Most permissive mode.  Tolerates heuristic regions and is
useful for exploratory diagnostics, never for a "proof of safe" claim. \\
\bottomrule
\end{tabular}
\end{center}

The mode affects the top-level verdict, not the set of reported bugs.  A real
shape bug is refuted identically in all three modes, which prevents users from
accidentally hiding an error by choosing a more conservative contract.

\subsection{Diagnostics, explanations, and adoption hooks}

The verifier produces source-mapped diagnostics with category, confidence, line
location, counterexample facts, and a suggested repair where the fix is
mechanical.  The \texttt{--explain} path reconstructs the inference chain that
led to a bug.  The \texttt{@tensorguard.checked} decorator and the upstream
\texttt{attach\_verifier} reference hook demonstrate an opt-in PyTorch
integration: a module is verified once at the first forward boundary, rejected
before the kernel runs if unsafe, or allowed transparently if safe.

\subsection{Security and safe loading}

Model verification often starts from untrusted Python files.  The repository
therefore contains a threat model, security policy, and a safe AST/fx-only path
that avoids dynamic import or \texttt{exec} when analyzing source.  Dynamic
execution remains available only in explicitly-labeled harnesses whose purpose
is to compare against real PyTorch behavior.  The production verifier treats
unsupported or opaque regions as \unknown{} in sound mode rather than executing
them.

\section{Program model and contracts}

\subsection{Verifiable fragment}

\path{src/verifiable_fragment.py} is the executable source of truth for what
\tg{} can analyze.  It defines a grammar for supported \texttt{nn.Module} and
\texttt{forward} constructs, tables of supported layers, tensor methods,
\texttt{torch.*} functions, and \texttt{torch.nn.functional} calls, plus an
unsupported-construct taxonomy.  The generated document
\path{VERIFIABLE_FRAGMENT.md} states the policy in user language:
unsupported code must be surfaced as \unknown{} in sound mode, never silently
passed.

Static checks catch data-dependent control flow, data-dependent iteration,
dynamic assertions, and tensor-to-scalar coercions such as \texttt{.item()}.
Tests run all statically detectable out-of-fragment categories through the
live verifier and confirm they yield \unknown{} under the sound contract.

\subsection{Soundness contract}

\path{src/soundness_contract.py} and \path{SOUNDNESS_CONTRACT.md} state two
directions separately:

\begin{description}[leftmargin=2em]
  \item[Refutation soundness.] A reported bug is backed by a concrete violated
  obligation.  The user should not have to distinguish a real bug from a solver
  artifact.
  \item[Verification soundness.] A \safe{} verdict in sound mode means no
  modeled bug class is missed within the verifiable fragment and the current
  operator confidence envelope.
\end{description}

The contract classifies constructs as \texttt{sound},
\texttt{over\_approximated}, \texttt{under\_approximated}, or \texttt{skipped}.
Known gaps are named rather than hidden.  For example, permissive modes may
return \safe{} on out-of-fragment code; sound mode abstains.  CEGAR's historic
SAFE-on-infeasible gap is closed by returning abstention rather than a false
\safe{}.

\subsection{Abstract domains}

The analysis tracks a reduced product of tensor facts.  The original shape
domain contains rank, concrete and symbolic dimensions, broadcast facts,
divisibility obligations, stride/layout, and permutation facts.  Additional
domains track dtype promotion, device placement, train/eval phase, and gradient
flow.  Reductions exchange facts across domains: a dtype cast into a floating
layer matters for shape-level operator validity; a device transfer changes the
device obligation at downstream ops; a branch guarded by training mode exposes
phase-dependent shape paths.

The domain-ablation artifacts distinguish verification domains from diagnostic
domains.  Shape, dtype, device, and gradient each earn their place by catching
bugs in a labeled corpus; phase currently records train/eval structure and
enables phase-specific diagnostics but is reported honestly where it does not
independently refute.

\subsection{Operator confidence}

Every registered operator transfer function has a machine-readable confidence
tag in \path{operator_confidence_table.json}.  The current table contains 117
operators: 75 \texttt{complete}, 36 \texttt{sound}, and 6 \texttt{heuristic}.
Unknown operators default to \texttt{heuristic}.  The CLI command
\texttt{tensorguard operator-confidence} surfaces these tags to users, and
sound mode folds heuristic operators into \unknown{} rather than claiming proof.

\section{Extraction frontends}

\tg{} offers multiple frontends that lower user code to a shared internal
representation.

\paragraph{AST frontend.}
The safe default parses source without importing the target module.  It reads
constructor assignments, \texttt{forward} dataflow, helper calls, subclass
patterns, literal \texttt{torch.randn} shapes, and simple input annotations.  It
is the right path for CI, untrusted files, and repository-wide scans.

\paragraph{\texttt{torch.fx} frontend.}
For live modules, the fx frontend traces idiomatic architectures and lowers
operator calls into the same transfer-function registry.  The trace-success
study covers 21 torchvision models, all traced and lowered, for 5,238 total
steps with 7.22\% unsupported operations recorded rather than hidden.

\paragraph{\texttt{torch.export}/Dynamo frontend.}
The export frontend reconciles PyTorch's capture path with the fx path.  The
current reconciliation artifact contains seven representative models, with both
frontends correct on all seven and no divergent verdicts.  Dynamic-shape export
constraints and export failure modes are surfaced through the same
verifiable-fragment taxonomy.

\paragraph{Second framework: Flax/JAX.}
The verifier core is framework-independent enough to support a Flax/JAX
frontend for \texttt{nnx.Sequential}-style models.  Linear, LayerNorm,
BatchNorm, and Dropout lower into the shared IR, shape bugs refute, and
unsupported layouts abstain soundly.  This demonstrates that the project is a
tensor-program verifier rather than a PyTorch-only parser trick.

\paragraph{Community stubs and autogeneration.}
Third-party layers are handled through governed declarative stubs: no executable
code, mandatory provenance, and executed conformance cases.  A stub autogenerator
derives safe shape stubs from real \texttt{torch.nn} constructor signatures when
the output contract is exactly known, and classifies uncertain layers as
unsupported rather than guessing.

\section{Constraint solving and refinement}

\subsection{Transfer functions}

Operator transfer functions are the heart of the analyzer.  They cover linear
layers, convolution and pooling families, normalization layers, embeddings,
matmul and batched matmul, broadcasting arithmetic, reshape/view/flatten,
transpose/permute, concatenation, split/chunk, gather/scatter, reductions,
attention blocks, dtype casts, device moves, train/eval-sensitive ops, and
gradient-flow constructs such as \texttt{detach}.  Frequency-weighted coverage
on a 13-model torchvision corpus is 94.86\%: 2,437 of 2,569 observed operator
occurrences are covered, with the remaining 132 occurrences listed by name.

\subsection{Solver obligations}

The verifier builds verification conditions over integer shape arithmetic,
finite-domain device and phase facts, dtype lattice facts, and gradient-flow
bits.  Z3 is the default backend because the implementation relies on its
incremental solver and user-propagator support.  A cvc5 comparison harness
checks a curated set of shape, device, phase, and reshape obligations through
both solvers and records concordant SAT/UNSAT answers, including on the
nonlinear reshape fragment.

The solver-domain flags
\texttt{--no-device-check}, \texttt{--no-phase-check}, and
\texttt{--no-grad-check} gate constraint generation itself, not merely a
post-hoc filter.  When a domain is disabled, its encoder emits no constraints
and its theory-specific checks are skipped; committed examples prove that each
flag can flip a verdict on real source.

\subsection{CEGAR predicate discovery}

The CEGAR loop discovers implicit shape contracts by accumulating predicates
from counterexamples.  If predicates discovered across refinement rounds are
jointly unsatisfiable, the conflict is promoted to a real bug of category
\texttt{cegar\_refined\_contract}.  This makes the iteration budget observable:
\texttt{--cegar-iterations 0} can produce no contract-level bug while a positive
budget surfaces the infeasible contract.  A convergence study proves the loop
terminates within a tight bound, and a Lean proof models the monotone predicate
accumulation argument.

\subsection{Specialized library contracts}

Seven high-value library checkers demonstrate the pattern for precision without
unsoundness:

\begin{itemize}[leftmargin=2em]
  \item \textbf{einops.}  \path{src/einops_verify.py} models
  \texttt{rearrange}, \texttt{reduce}, and \texttt{repeat}; it catches
  non-divisible decompositions, underdetermined splits, axis-set mismatches,
  duplicate axes, anonymous-axis errors, rank mismatches, and ellipsis cases.
  A 3,000-case differential fuzz battery checks verdicts and output shapes
  against the real \texttt{einops} package.  Symbolic divisibility that cannot
  be decided statically is not refuted.
  \item \textbf{Scaled dot-product attention.}  \path{src/sdpa_verify.py}
  models the rank, head-dimension, batch/head broadcasting, and mask
  broadcasting contract of \texttt{F.scaled\_dot\_product\_attention}.  A
  400-case differential battery checks it against real PyTorch.  Backend-specific
  key/value sequence-length behavior is not hard-refuted where PyTorch's default
  fused backend accepts it.
  \item \textbf{\texttt{torch.distributions}.}
  \path{src/distributions_verify.py} resolves batch shapes, event shapes, and
  \texttt{log\_prob} output shapes for scalar-event distributions,
  \texttt{Categorical}, \texttt{MultivariateNormal}, and
  \texttt{Independent}.  The contract is intentionally over usable
  \texttt{sample}/\texttt{log\_prob} behavior, not just constructors, because
  PyTorch can build a malformed multivariate distribution that only fails when
  used.  Differential tests compare against real PyTorch distributions with
  support validation disabled so the oracle isolates shape failures.
  \item \textbf{Complex views and FFT.}
  \path{src/complex_verify.py} models \texttt{torch.view\_as\_real},
  \texttt{torch.view\_as\_complex}, and the core \texttt{torch.fft} family.
  It captures PyTorch's concrete shape rules for \texttt{n}, \texttt{s}, and
  \texttt{dim}, including no-op \texttt{fftn(..., dim=())}, Hermitian
  half-spectrum lengths, inverse real-FFT length recovery, out-of-range and
  duplicate dimensions, and the exact dtype edge cases: \texttt{rfft} rejects
  complex inputs, \texttt{view\_as\_complex} accepts only half/float/double
  real views, and FFT kernels reject half, bfloat16, and complex-half where the
  live dispatcher does.  Unknown dtypes and symbolic lengths abstain instead of
  refuting.  \path{tests/test_complex_verify.py} differentially checks the rules
  against real PyTorch across view, single-axis FFT, and multi-axis FFT cases.
  \item \textbf{PyTorch named tensors.}
  \path{src/named_tensor_verify.py} models \texttt{refine\_names} and
  \texttt{align\_to}: an existing concrete axis name cannot be renamed or
  demoted, Ellipsis carries omitted axes in torch's order, new target names
  insert singleton dimensions, and unnamed dimensions are moved only when
  covered by Ellipsis.  A conservative source checker reports invalid literal
  named-tensor calls in real code.  Differential tests compare random
  refine/alignment cases against real PyTorch named tensors, including
  duplicate names, invalid identifiers, explicit \texttt{None} axes, and
  symbolic sizes carried through unchanged.
  \item \textbf{\texttt{torch.vmap}/functorch.}
  \path{src/vmap_verify.py} models the higher-order shape transfer used by
  \texttt{torch.vmap} and \texttt{torch.func.vmap}: mapped input axes are
  removed before the body, concrete mapped batch sizes must agree, and the
  shared batch axis is reinserted according to \texttt{out\_dims}.  The checker
  supports nested output pytrees and treats symbolic batch dimensions by
  abstention.  For PyTorch's version-sensitive constant-output corner, it
  records an unknown shape rather than refuting.  A 500-case randomized
  differential test checks mapped-axis removal and out-dim insertion against
  real \texttt{vmap}.
  \item \textbf{\texttt{torch.sparse}.}
  \path{src/sparse_verify.py} models COO sparse/dense rank agreement and the
  validated CSR/CSC/BSR/BSC compressed-layout contract: batch metadata, nnz
  agreement, compressed-axis lengths, positive block sizes, and block
  divisibility.  It also flags compressed tensors whose values' dense tail
  disagrees with the requested tensor shape: PyTorch may construct these
  metadata-inconsistent tensors, but \texttt{to\_dense()} fails, so \tg{}
  reports them as unusable layouts.  The tests compare valid and invalid cases
  against real PyTorch constructors with \texttt{check\_invariants=True}, plus
  the live dense-materialization failure and symbolic no-false-positive probes.
  \item \textbf{\texttt{grid\_sample}/\texttt{affine\_grid}.}
  \path{src/grid_sample_verify.py} models the 2-D and 3-D geometric sampler
  contracts: input/grid rank, batch equality, coordinate-grid trailing size,
  non-empty input spatial axes, bicubic's 4-D restriction, dtype agreement, and
  \texttt{affine\_grid}'s theta-matrix and positive-size requirements.  The
  tests replay hand edge cases and randomized cases against real PyTorch,
  including the easy-to-get-wrong distinction that \texttt{grid\_sample} accepts
  empty output grids while rejecting empty input spatial axes.
\end{itemize}

\section{Beyond \texttt{forward}: deployment and training hazards}

The repository extends static checking beyond architecture shape.

\paragraph{Training-loop hazards.}
\path{src/training_loop_checks.py} analyzes optimizer and gradient-flow
patterns: detached losses, missing \texttt{zero\_grad}, missing
\texttt{optimizer.step}, \texttt{zero\_grad} after \texttt{backward}, and fp16
autocast without a \texttt{GradScaler}.  The reproducibility harness runs
equivalent seeded PyTorch loops and confirms the static verdict matches the
runtime symptom.

\paragraph{Distributed tensor parallelism.}
\path{src/tensor_parallel_checks.py} models Megatron-style column-parallel and
row-parallel linear stacks.  It checks shard divisibility, contracted-dimension
agreement, and communication flag consistency.  The harness hand-shards a
reference stack across simulated ranks and verifies that correct configs match
the unsharded forward while inconsistent configs are both statically flagged
and fail under actual sharding.

\paragraph{Quantization and export.}
\path{src/quant_export_checks.py} catches asymmetric quant/dequant boundaries,
quantized tensor arithmetic that should use \texttt{FloatFunctional}, and
export-unsafe data-dependent constructs.  The live PyTorch harness confirms
the quantized-add hazard raises the expected runtime exception and that export
failures line up with the verifiable-fragment grammar.

\paragraph{ONNX, \texttt{torch.compile}, and packaging gates.}
The repository includes ONNX hardening, post-export \texttt{check\_model}
assertions, and torch.export/AOTInductor packaging parity.  These gates are
framed as deployment checks: they do not broaden the soundness contract unless
their corresponding transfer rules and runtime conformance tests are present.

\section{Formal assurance}

\subsection{Lean proof corpus}

The Lean 4 development under \path{lean/TensorGuard/} mechanizes the proof
obligations that are most dangerous to leave as prose: operator rules,
reduced-product algebra, CEGAR bounds, fragment modes, device/dtype/phase/grad
transfers, shape-chain transfers, and SMT encoding faithfulness.  The whole
library audit test reads the root import graph, discovers every public theorem
and lemma, generates \texttt{\#print axioms} commands for each, elaborates them
through the Lean kernel, and asserts that no declaration depends on
\texttt{sorryAx} and that every reported axiom is one of
\texttt{propext}, \texttt{Classical.choice}, or \texttt{Quot.sound}.  The
current discovered count is 371 public theorems.

The proof corpus is paired with Python tests rather than left disconnected.
Conformance tests map Lean theorem names to randomized property tests of the
Python code.  Handler parity checks compare Lean-extracted or Lean-modeled
rules with the Python transfer functions and, where applicable, the live torch
dispatcher.  This is important: a theorem about a toy rule is not enough unless
the implementation that users run is pinned to the same semantics.

\subsection{Encoding faithfulness}

Several Lean modules address the trust boundary between abstract transfer
rules and solver encodings.  \path{lean/TensorGuard/SmtEncoding.lean} proves that the generated
Z3 obligations for dtype/device/phase/gradient facts faithfully represent the
Lean semantics.  \path{lean/TensorGuard/CrossDomain.lean} lifts shape/device interactions.
Additional modules cover dtype promotion chains, gradient-flow chains,
device-placement chains, phase chains, broadcast chains, rank-transfer chains,
flatten, cat, embedding, reshape \texttt{-1} inference, Linear, Conv1d/Conv2d,
ConvTranspose1d, pooling, LayerNorm, PixelShuffle, AdaptiveAvgPool2d,
Unflatten, and BatchNorm.

\subsection{Soundness footprint}

The paper-grade soundness story is deliberately scoped.  The repository
distinguishes Lean-audited handlers, pen-and-paper handlers, tested-only
handlers, and uncovered handlers.  Per-block soundness-footprint artifacts pin
real refutations to the concrete handler set they traversed, including the
specific Lean theorem or rule that discharges each audited handler.  This
prevents a misleading global statement such as "Lean proves the whole tool" from
masking uncovered long-tail operators.

\section{Benchmark corpora}

The benchmark strategy uses multiple corpora because no single corpus can answer
every reviewer question.

\begin{longtable}{p{0.24\linewidth} p{0.14\linewidth} p{0.52\linewidth}}
\caption{Committed corpora and what each one proves.}\label{tab:corpora}\\
\toprule
\textbf{Corpus / source} & \textbf{Size} & \textbf{Purpose} \\
\midrule
\endfirsthead
\toprule
\textbf{Corpus / source} & \textbf{Size} & \textbf{Purpose} \\
\midrule
\endhead
\path{real_benchmarks/} & 16 & Frozen ground-truth real-code corpus: 8 clean
and 8 buggy modules across shape, device, phase, and gradient domains, each
content-addressed by SHA-256.  Six of eight buggy modules raise a real eager
PyTorch RuntimeError on this host; the CUDA-only and silent-gradient cases are
documented. \\
\path{experiments_v5/github_bug_mining/} & 2,704 & Public GitHub issue/PR
snapshot mined from verbatim PyTorch runtime error signatures: 2,394 shape and
310 device bugs across eight categories. \\
\path{corpus_extended/} & 227 & Parameterized, redistributable generated
corpus: 153 buggy and 74 clean cases, all labels validated by real PyTorch
execution and content-addressed. \\
Blind held-out split & 186 & Pre-registered split disjoint from development
parameter grids: 138 buggy and 48 clean cases, scored exactly once. \\
False-positive stress corpus & 101 & Clean models built from tricky but valid
shape, broadcast, reshape, concat, and normalization patterns. \\
Natural-distribution study & 29 & Clean public-style architectures across
fourteen families; measures definite-answer coverage on ordinary code. \\
Differential dispatcher corpus & 2,000 & Random modules judged by both eager
PyTorch and \tg{} sound mode; tests exact agreement with the live dispatcher. \\
Hypothesis module-AST corpus & 800 & Structured generated modules with
shrinking, testing the full-module soundness property rather than isolated
operators. \\
Mutation corpus & 376 genuine mutants & Mutations of validated-clean parents;
only mutants that genuinely raise under eager PyTorch are admitted. \\
\bottomrule
\end{longtable}

\paragraph{Why not only real bugs?}
The frozen real corpus provides provenance; the GitHub-mined snapshot provides
scale; generated corpora provide redistributability and controlled labels;
blind splits test overfitting; fuzzing and mutation testing search for
regressions in both directions.  The combination is what makes the evaluation
credible.

\section{Main empirical results}

\subsection{Frozen real-code benchmark}

Table~\ref{tab:confusion} reports the committed confusion matrices on the
16-case frozen corpus.  Baselines are actual tools or executable procedures:
a no-op predictor, a runtime forward smoke test, a runtime forward+backward
test that also inspects gradients, and PyTea built from source and run live on
generated entry wrappers.

\begin{table}[h]
\centering
\small
\begin{tabular}{lrrrrrr}
\toprule
\textbf{Method} & \textbf{TP} & \textbf{FP} & \textbf{TN} & \textbf{FN}
& \textbf{Precision} & \textbf{Recall} \\
\midrule
\tg{} & 8 & 0 & 8 & 0 & 1.000 & 1.000 \\
runtime backward & 8 & 0 & 8 & 0 & 1.000 & 1.000 \\
runtime forward & 7 & 0 & 8 & 1 & 1.000 & 0.875 \\
PyTea & 6 & 0 & 8 & 2 & 1.000 & 0.750 \\
no-op & 0 & 0 & 8 & 8 & -- & 0.000 \\
\bottomrule
\end{tabular}
\caption{Frozen 16-case corpus; the backing artifact is evaluation/confusion\_matrices.json.}
\label{tab:confusion}
\end{table}

The runtime backward baseline ties on this small corpus because each bug is
exercised by the seeded run and because the gradient-detach case can be caught
by checking parameter gradients.  That tie is useful: it shows the comparison is
not rigged against dynamic tools.  The static advantage appears on latent bugs
where one concrete execution cannot exercise the failing path.

\subsection{Latent bugs}

The hard-recall corpus contains eight bugs that are all proven genuine but
silent under a single seeded train-mode forward/backward baseline: three
eval-only phase bugs, three path-flag bugs, and two silent parameter-freeze
bugs.  The strongest dynamic baseline catches 0/8.  \tg{} catches 6/8,
covering the phase and path families in full; the two misses are explicitly
tagged as post-construction \texttt{requires\_grad=False} freezes outside the
current gradient-domain model.  This is exactly the failure mode static analysis
should improve: not "faster execution", but detection of bugs no chosen
execution observes.

\subsection{False-positive and false-negative hunts}

\begin{table}[h]
\centering
\small
\begin{tabular}{p{0.34\linewidth} p{0.18\linewidth} p{0.36\linewidth}}
\toprule
\textbf{Study} & \textbf{Population} & \textbf{Result} \\
\midrule
Sound-mode clean hunt & 80 clean models & 0 false positives; 80/80 \safe{};
no abstentions. \\
Differential clean fuzzing & 200 random clean models & 0 false positives;
200/200 \safe{}; ranks split 94 rank-2 and 106 rank-4. \\
Negative fuzzing & 281 genuine injected faults & 281/281 caught; 0 false
negatives across Linear-in, Linear-out, Conv-in, and bad-reshape families. \\
Disagreement triage & 481 fuzzed clean/faulty models & 0 TensorGuard/runtime
disagreements; 50 minimal bug reproducers frozen as regressions. \\
Extended corpus & 227 cases & 153/153 runtime-failing cases caught; 0 false
positives on 74 clean cases in both balanced and sound modes. \\
Blind split & 186 cases & 138/138 held-out buggy cases caught; 0 false
positives on 48 clean cases; zero overfitting gap. \\
Live dispatcher differential & 2,000 modules & Exact agreement with eager
PyTorch: zero soundness violations, zero false alarms, zero abstentions. \\
Mutation testing & 376 genuine mutants & Sound mode kills every admitted
mutant; no genuine bug reported \safe{}. \\
\bottomrule
\end{tabular}
\caption{Independent robustness studies.  Each row is generated by a separate
script and pinned by tests.}
\label{tab:robustness}
\end{table}

The most important engineering property is not a single high recall number; it
is the absence of easy counterexamples.  The project therefore attacks the tool
from both sides: generate clean models and demand no false alarms, then inject
real faults and demand no misses.  When these hunts find no natural
disagreements, the repository still freezes a regression suite of minimal
buggy/clean siblings so future changes cannot accidentally regress.

\subsection{Historical and mined bugs}

The historical 60-bug corpus reports two clearly-labeled regimes.  In the
default headline regime, \tg{} produces \rp{} verdicts on 53/60 bugs at the
high-confidence threshold, with seven silent misses.  In the raw-refute regime,
which lifts per-bug input shapes and allows CEGAR depth three, it refutes 56/60,
55 of them high-confidence.  The audit harness recomputes both numbers from
\path{reproducibility/reproduce_headline_60bug.json}.  The larger GitHub-mined
snapshot is not a precision/recall benchmark by itself; it is a frozen,
content-addressed source of real-world bug signatures and categories used to
drive future corpus growth.

\subsection{Operator and frontend coverage}

Coverage is measured two ways.  The registry-level table tags 117 transfer
functions by confidence.  The real-model frequency study asks how often those
rules appear on 13 torchvision architectures: 94.86\% of 2,569 operator
occurrences are covered.  The fx trace-success study asks whether the frontend
can capture real models: 21/21 torchvision models trace and lower, totaling
5,238 steps; unsupported operations are counted and reported, not ignored.

\subsection{Localization and developer effort}

For refuted real bugs with author-placed \texttt{\# BUG} markers, the
localization proxy measures how many lines a developer inspects to reach the
true bug.  Median assisted effort is two lines; median unaided effort is 7.5
lines; the median reduction factor is 3.0 with a bootstrap interval of
[1.33125, 5.875].  \tg{} helps on 12 of 14 bugs and hurts on two, which are kept
in the artifact.  The repository also contains a pre-registered human-study
protocol, because this proxy is not claimed to be a human RCT.

\subsection{Statistical testing}

\path{evaluation/significance.py} runs paired exact McNemar tests with
Holm-Bonferroni correction and paired bootstrap confidence intervals over the
confusion-matrix predictions.  On the 16-case frozen corpus, \tg{} significantly
beats the no-op floor after correction, never loses a discordant pair to any
baseline, and does not overclaim significance over PyTea at this sample size.
\path{reproducibility/effect_sizes.py} adds paired effect sizes and both FWER
and FDR corrections.  \path{reproducibility/statistical_power.py} reports
sample-size justifications and rule-of-three bounds for zero-false-alarm claims.

\subsection{Scaling}

The scaling study sweeps feed-forward depth from one to 64 layers.  The
structural-work metric is linear in depth; the wall-clock companion has a
log-log scaling exponent of 1.69 and $R^2=0.988$, below cubic and far from an
exponential blow-up.  Median wall time grows from about 12 ms at depth one to
about 4.26 s at depth 64 on the recorded host.

\section{Reproducibility machinery}

\subsection{One command for the evidence base}

The \path{Makefile} exposes \texttt{make reproduce}, \texttt{make
reproduce-check}, \texttt{make reproduce-full}, and task-specific targets such
as \texttt{make precision-recall}, \texttt{make sound-fp},
\texttt{make diff-fuzz}, \texttt{make neg-fuzz}, and \texttt{make
paper-evidence}.  The orchestrator \path{reproducibility/reproduce_all.py}
regenerates deterministic generated docs, tables, corpus manifests, benchmark
artifacts, and the headline 60-bug figure before running the numeric-claim
audit.

\subsection{Numeric-claim audit}

Every headline numeric claim in \path{README.md}, the conference paper, and the
workshop paper is registered in \path{reproducibility/audit_numeric_claims.py}.
For each claim, the harness checks that the number still appears in the cited
source and recomputes it from the backing artifact.  Claims can be
\texttt{VERIFIED}, \texttt{MISMATCH}, \texttt{QUALIFIED\_REGIME},
\texttt{QUALIFIED\_ENV}, \texttt{ORPHAN}, or \texttt{SOURCE\_MISSING}.  The
current committed audit has no uncovered README numeric tokens.

\subsection{Artifact-evaluation capsule}

The repository includes an artifact-evaluation package under
\path{docs/artifact/} and a reproducibility capsule with pinned wheel locks,
a Dockerfile, and \path{capsule/reproduce.sh}.  The capsule verifies the live
interpreter, regenerates deterministic artifacts, checks byte identity, and
reruns the numeric audit.  Badge evidence is mapped to in-tree paths for
Available, Functional, Reusable, and Results Reproduced criteria.

\subsection{Dashboard and release gates}

\path{docs/dashboard/index.html} is a static site generated entirely from
committed JSON artifacts.  \path{evaluation/dashboard_baseline.json} freezes
quality and integrity metrics so regressions are reviewable diffs rather than
silent changes.  \path{scripts/benchmark_delta.py} compares release snapshots
with metric-direction awareness: recall must not drop, false positives must not
rise, integrity populations must not shrink, and machine-dependent wall-clock
metrics are reported without gating.  Coverage gates hold selected public API
and integration modules above 90\% line coverage.

\section{Deployment and governance}

\subsection{Packaging}

\tg{} is packaged as a typed Python project with \texttt{py.typed}, a pinned
\texttt{z3-solver} dependency range, a CLI entrypoint, a pre-commit hook, and a
pytest plugin.  The roadmap includes PyPI, conda-forge, Docker, and pipx paths;
the repository already contains the policy and harness pieces needed to make
those releases reproducible.

\subsection{Community benchmark and stubs}

The leaderboard scores tools on the frozen real benchmark corpus by recomputing
metrics from raw verdict files.  External tools cannot self-report inflated
numbers.  The community stub registry gives framework and library authors a
safe path to extend coverage: manifests must be declarative, carry provenance,
and ship executed conformance cases.  The governance docs make reviewability
part of the technical contract.

\subsection{Documentation and tutorials}

The README intentionally surfaces limitations near the feature list.  The
Getting Started guide, "what it cannot do yet" document, Colab-style tutorials,
examples gallery, artifact appendix, user-study protocol, PyTorch RFC draft,
growth playbook, and dashboard are all tested for path existence or executable
snippets where appropriate.  Documentation is treated as a checked artifact,
not post-hoc marketing prose.

\section{Threats to validity and limitations}

\paragraph{Not all PyTorch is in scope.}
PyTorch is too large for a single static checker to model completely.  The
verifiable fragment is explicit, and unsupported constructs produce \unknown{}
in sound mode.  This means \safe{} is meaningful, but \unknown{} is a normal and
expected answer on dynamic Python.

\paragraph{Some domains are diagnostic or partial.}
Phase facts are useful for train/eval diagnostics, but the current marginal
contribution study classifies phase as diagnostic-only where it does not
independently refute.  The gradient domain catches \texttt{detach}-style
forward-flow bugs but currently misses post-construction parameter freezes such
as \texttt{requires\_grad=False}; those misses are tagged in the artifact.

\paragraph{Generated corpora are not natural distributions.}
Generated and mutated corpora are valuable because they are redistributable and
have executable labels, but they do not replace naturally-occurring bugs.  The
project therefore reports them alongside real benchmarks, GitHub-mined issue
snapshots, and natural-distribution clean samples.

\paragraph{Some regeneration is environment-qualified.}
Lean, CUDA, and HuggingFace model downloads are not required for the standard
CPU-only reproduction path.  Claims depending on those environments are
recorded with regeneration commands and classified as environment-qualified by
the audit.

\paragraph{Formal proof scope is explicit.}
Lean proves many central rules and audits the public theorem corpus, but the
entire Python implementation, parser, and every long-tail handler are not all
inside a single theorem.  The soundness-footprint artifacts state which verdicts
traverse Lean-audited, pen-and-paper, tested-only, or uncovered handlers.

\section{Related tools}

\paragraph{Runtime shape libraries.}
Libraries such as jaxtyping and torchtyping check tensor contracts at runtime
and often require annotations.  They are useful when the annotated path is
executed, but they do not statically prove unexecuted branches and do not cover
device, phase, export, or distributed hazards by default.

\paragraph{Static tensor analyzers.}
PyTea is the closest academic baseline for static PyTorch shape checking.  The
repository builds and runs PyTea live where possible, reports a shape-only
subcorpus where tools tie, and separately reports where \tg{}'s additional
domains catch bugs PyTea is not designed to model.

\paragraph{PyTorch tracing and export.}
\texttt{torch.fx}, \texttt{torch.export}, TorchScript, and torch.compile can
surface shape failures by tracing or compiling with concrete examples.  They are
powerful dynamic or capture-time tools, but they require instantiation and
example inputs.  \tg{} uses these paths as frontends and baselines while
retaining an input-free static path.

\paragraph{General static analyzers.}
Mypy, Pyright, and similar type checkers catch Python type errors but do not
reason about tensor dimensions or cross-device placement.  \tg{} is
complementary: it focuses on tensor metadata contracts rather than replacing
Python type checking.

\section{What makes the repository unusually strong}

A typical research prototype has one impressive idea and a thin evaluation.
\tg{} has a different shape: its claim is that static tensor verification can be
made operationally trustworthy.  That requires several reinforcing layers:

\begin{enumerate}[leftmargin=2em]
  \item a strict contract that refuses to overclaim;
  \item transfer functions tagged by confidence;
  \item real and generated corpora with executable labels;
  \item fuzzing, mutation, dispatcher differential tests, and regression
  minimization;
  \item formal proofs for the rules most likely to cause unsoundness;
  \item claim audits so prose cannot drift from artifacts;
  \item CI gates that make benchmark regressions reviewable;
  \item adoption surfaces that integrate with how developers actually ship
  PyTorch code.
\end{enumerate}

This combination is what makes the project a plausible candidate for upstream
PyTorch influence: not that it already models every operator, but that it
provides a disciplined path for growing coverage without sacrificing trust.

\section{Next research and adoption frontier}

The local roadmap now points beyond the completed campaign toward a
best-paper-award and 1,000-star release, with one hundred concrete future work
items numbered after the completed campaign.  The most important next
directions are not more marketing metrics; they are higher-value coverage and
stronger evidence:

\begin{itemize}[leftmargin=2em]
  \item deeper operator coverage for MultiheadAttention, split/chunk, masked
  selection, nonzero, and exact einsum diagonal semantics;
  \item deployment correctness for quantization, mixed precision, dynamic
  \texttt{torch.export} constraints, AOTInductor layouts, ONNX opsets, GGUF
  export, FSDP2, and pipeline-parallel stage boundaries;
  \item proof depth for einops, SDPA, chunk/split, distributions, named-tensor
  alignment, grid sampling, contiguity/stride algebra, and whole-module subject
  reduction over the expanded operator set;
  \item evaluation scale: 1,000+ provenance-rich real bugs, stratified recall
  and false-positive confidence intervals, registered held-out protocols,
  head-to-head baselines, and PR-history survival analysis;
  \item adoption: hosted demo, one-line GitHub Action and pre-commit setup,
  editor clients, Code Scanning trends, web explanations, Lightning/HF Trainer
  adapters, model-hub badges, a real-model gallery, and a public leaderboard
  refresh cadence.
\end{itemize}

\section{Conclusion}

\tg{} turns tensor metadata checking from a runtime accident into an auditable
static contract.  It is already useful as a developer tool: it catches shape,
device, dtype, phase, gradient, training-loop, distributed, quantization,
export, einops, SDPA, distribution, complex FFT/view, named-axis, vmap,
sparse-layout, and sampler-grid hazards
before execution; it reports precise
diagnostics; it integrates with CI and code scanning; and it defaults to
abstention rather than false certainty when the program leaves its fragment.
It is also unusually well-supported as a research artifact: every headline
number is tied to a regenerable artifact, every corpus has provenance, every
known limitation is surfaced, and the formal proof story is executable enough
to fail CI when it drifts.

\newpage
\appendix

\section{Evidence index}

\begin{longtable}{p{0.27\linewidth} p{0.33\linewidth} p{0.30\linewidth}}
\caption{Selected repository artifacts backing the paper.}\\
\toprule
\textbf{Question} & \textbf{Artifact / command} & \textbf{Tests / gate} \\
\midrule
\endfirsthead
\toprule
\textbf{Question} & \textbf{Artifact / command} & \textbf{Tests / gate} \\
\midrule
\endhead
What programs are soundly verified? &
\path{SOUNDNESS_CONTRACT.md}, \path{src/soundness_contract.py} &
\path{tests/test_soundness_contract.py},
\path{tests/test_soundness_mode.py} \\
What source fragment is supported? &
\path{VERIFIABLE_FRAGMENT.md}, \path{src/verifiable_fragment.py} &
\path{tests/test_verifiable_fragment_spec.py} \\
Which operators are complete, sound, or heuristic? &
\path{operator_confidence_table.json},
\texttt{tensorguard operator-confidence} &
\path{tests/test_operator_confidence.py} \\
Does disabling a domain skip solver work? &
\path{reproducibility/check_flag_demo.json} &
\path{tests/test_solver_domain_gating.py} \\
Does CEGAR change real bug results? &
\path{tests/test_cegar_refined_contract.py},
\path{reproducibility/cegar_depth_ablation.md} &
\path{tests/test_cegar_depth_ablation.py} \\
How does \tg{} compare to baselines? &
\path{evaluation/confusion_matrices.json},
\path{evaluation/significance.json} &
\path{tests/test_precision_recall.py},
\path{tests/test_significance.py} \\
Does sound mode false-alarm on clean models? &
\path{evaluation/sound_mode_fp.json} &
\path{tests/test_sound_mode_fp.py} \\
Does random clean fuzzing find false positives? &
\path{evaluation/diff_fuzz.json} &
\path{tests/test_diff_fuzz.py} \\
Does injected-fault fuzzing find false negatives? &
\path{evaluation/neg_fuzz.json} &
\path{tests/test_neg_fuzz.py} \\
Are disagreements minimized and frozen? &
\path{evaluation/minimize_demo.json},
\path{evaluation/triage_regressions.json} &
\path{tests/test_minimize.py}, \path{tests/test_triage.py} \\
Is there an extended redistributable corpus? &
\path{corpus_extended/manifest.json},
\path{reproducibility/corpus_extended_score.md} &
\path{tests/test_corpus_extended.py} \\
Is there a blind split? &
\path{corpus_extended/PRE_REGISTRATION.md},
\path{reproducibility/blind_split_eval.md} &
\path{tests/test_blind_split.py} \\
Does it agree with the live torch dispatcher? &
\path{reproducibility/differential_dispatcher.md} &
\path{tests/test_differential_dispatcher.py} \\
Does mutation testing catch genuine bugs? &
\path{reproducibility/mutation_clean_models.md} &
\path{tests/test_mutation_clean_models.py} \\
Does the localizer reduce effort? &
\path{evaluation/localization_effort.json} &
\path{tests/test_localization_effort.py} \\
How much operator coverage exists on real models? &
\path{evaluation/operator_frequency.json},
\path{evaluation/fx_trace_success.json} &
\path{tests/test_operator_frequency.py},
\path{tests/test_fx_trace_success.py} \\
Are einops contracts faithful? &
\path{src/einops_verify.py},
\path{src/einops_source.py} &
\path{tests/test_einops_verify.py},
\path{tests/test_einops_source.py} \\
Is SDPA modeled faithfully? &
\path{src/sdpa_verify.py} &
\path{tests/test_sdpa_verify.py} \\
Are probabilistic batch/event shapes modeled faithfully? &
\path{src/distributions_verify.py} &
\path{tests/test_distributions_verify.py} \\
Are named-tensor refinements and alignments modeled faithfully? &
\path{src/named_tensor_verify.py} &
\path{tests/test_named_tensor_verify.py} \\
Is \texttt{vmap} batched-shape transfer modeled faithfully? &
\path{src/vmap_verify.py} &
\path{tests/test_vmap_verify.py} \\
Are sparse layout and blocksize contracts modeled faithfully? &
\path{src/sparse_verify.py} &
\path{tests/test_sparse_verify.py} \\
Are geometric sampler grids modeled faithfully? &
\path{src/grid_sample_verify.py} &
\path{tests/test_grid_sample_verify.py} \\
Are training-loop hazards checked? &
\path{reproducibility/training_loop_hazards.md} &
\path{tests/test_training_loop_checks.py} \\
Are tensor-parallel configs checked? &
\path{reproducibility/tensor_parallel_sharding.md} &
\path{tests/test_tensor_parallel_checks.py} \\
Are quant/export hazards checked? &
\path{reproducibility/quant_export_safety.md} &
\path{tests/test_quant_export_checks.py} \\
Is the Lean library sorry-free? &
\path{lean/TensorGuard/},
\path{tests/test_lean_whole_pipeline_audit.py} &
\texttt{lake build TensorGuard} in slow environments; axiom audit tests \\
Are headline numbers audited? &
\path{reproducibility/numeric_claims_audit.json} &
\path{tests/test_numeric_claims_audit.py} \\
Can the evidence be reproduced? &
\texttt{make reproduce}, \texttt{make reproduce-check},
\path{capsule/reproduce.sh} &
\path{tests/test_reproduce_harness.py},
\path{tests/test_capsule_manifest.py} \\
\bottomrule
\end{longtable}

\section{Experiment catalogue}

The repository deploys the following benchmark and validation families:

\begin{enumerate}[leftmargin=2em]
  \item \textbf{Ground-truth real benchmarks:} content-addressed clean and buggy
  \texttt{nn.Module}s with provenance and eager-runtime validation.
  \item \textbf{GitHub mining:} frozen public issue/PR snapshot from verbatim
  PyTorch runtime signatures.
  \item \textbf{Precision/recall baselines:} TensorGuard, PyTea, runtime
  forward, runtime backward, and no-op on the same corpus.
  \item \textbf{Statistical rigor:} McNemar tests, bootstrap intervals, effect
  sizes, multiple-comparison correction, and power analysis.
  \item \textbf{False-positive hunts:} clean generated templates, random
  architecture fuzzing, false-positive stress grids, and natural-distribution
  clean models.
  \item \textbf{False-negative hunts:} injected genuine faults, mutation
  testing, latent-bug corpora, and full dispatcher differential tests.
  \item \textbf{Frontend studies:} fx trace success, export/fx reconciliation,
  cross-version stability, cross-Python determinism, and safe blocked-import
  checks.
  \item \textbf{Operator studies:} registry coverage, frequency-weighted real
  model coverage, conformance oracle checks, shape algebra properties, einops
  fuzzing, SDPA fuzzing, and probabilistic distribution-shape fuzzing.
  \item \textbf{Deployment checks:} training-loop hazards, tensor-parallel
  sharding, quantization/export safety, ONNX/export package gates, and upstream
  hook demos.
  \item \textbf{Reproducibility checks:} numeric claims audit, artifact badges,
  capsule manifest, evidence dashboard, benchmark-delta gate, and one-command
  paper evidence generation.
\end{enumerate}

\section{Capability matrix}

This matrix is intentionally implementation-facing: it records what a user,
reviewer, or future contributor can find in the repository today and which part
of the system owns the behavior.

\begin{longtable}{p{0.20\linewidth} p{0.34\linewidth} p{0.34\linewidth}}
\caption{Major shipped capabilities grouped by surface area.}\\
\toprule
\textbf{Surface} & \textbf{Capability} & \textbf{Primary implementation or evidence} \\
\midrule
\endfirsthead
\toprule
\textbf{Surface} & \textbf{Capability} & \textbf{Primary implementation or evidence} \\
\midrule
\endhead
Core verifier & Zero-annotation \texttt{nn.Module} architecture verification
with shape, dtype, device, phase, gradient, stride, and layout facts &
\path{src/api.py}, \path{src/smt/solver.py} \\
Core verifier & Three-valued verdicts \safe{}/\unknown{}/\unsafe{} with
\texttt{sound}, \texttt{balanced}, and \texttt{heuristic} modes &
\path{src/soundness_contract.py}, \path{tests/test_soundness_mode.py} \\
Core verifier & Solver-domain gates for device, phase, and gradient checks that
remove constraints before Z3, not after the verdict &
\path{tests/test_solver_domain_gating.py},
\path{reproducibility/check_flag_demo.json} \\
Core verifier & CEGAR-discovered infeasible contracts promoted to real bugs &
\path{tests/test_cegar_refined_contract.py},
\path{reproducibility/cegar_depth_ablation.md} \\
Operator model & Confidence-tagged transfer registry and CLI inspection &
\path{src/operator_confidence.py},
\path{operator_confidence_table.json} \\
Operator model & Precise reshape/view/flatten, broadcasting, einsum,
convolution, normalization, dtype-promotion, and device-transfer semantics &
\path{tests/test_shape_algebra_properties.py},
\path{tests/test_lean_conformance.py} \\
Library model & einops \texttt{rearrange}/\texttt{reduce}/\texttt{repeat}
checker with source-level bug reports &
\path{src/einops_verify.py}, \path{src/einops_source.py} \\
Library model & SDPA shape and mask checker for attention hot paths &
\path{src/sdpa_verify.py}, \path{tests/test_sdpa_verify.py} \\
Library model & \texttt{torch.distributions} batch/event/log-prob shape checker &
\path{src/distributions_verify.py},
\path{tests/test_distributions_verify.py} \\
Library model & PyTorch named-tensor \texttt{refine\_names}/\texttt{align\_to}
checker with source-level bug reports &
\path{src/named_tensor_verify.py}, \path{tests/test_named_tensor_verify.py} \\
Library model & \texttt{torch.vmap}/functorch mapped-axis removal and
\texttt{out\_dims} insertion &
\path{src/vmap_verify.py}, \path{tests/test_vmap_verify.py} \\
Library model & \texttt{torch.sparse} COO/CSR/CSC/BSR/BSC layout and blocksize
checks &
\path{src/sparse_verify.py}, \path{tests/test_sparse_verify.py} \\
Library model & \texttt{grid\_sample}/\texttt{affine\_grid} sampler-rank,
batch, dtype, and output-grid contracts &
\path{src/grid_sample_verify.py}, \path{tests/test_grid_sample_verify.py} \\
Frontends & Safe AST source parser for untrusted model files &
\path{src/safe_loader.py}, \path{tests/test_security.py} \\
Frontends & \texttt{torch.fx} trace lowering across real torchvision models &
\path{evaluation/fx_trace_success.json} \\
Frontends & \texttt{torch.export}/Dynamo reconciliation with fx &
\path{evaluation/frontend_reconciliation.json} \\
Frontends & Flax/JAX frontend lowering into the shared IR &
\path{src/flax_extractor.py}, \path{tests/test_flax_frontend.py} \\
Extensibility & Governed declarative community stubs with provenance and
executed conformance cases &
\path{community_stubs/}, \path{src/stub_governance.py} \\
Training & Static checks for detached loss, missing or misordered
\texttt{zero\_grad}, missing \texttt{step}, and unsafe fp16 autocast &
\path{src/training_loop_checks.py},
\path{reproducibility/training_loop_hazards.md} \\
Distributed & Tensor-parallel shard divisibility and communication-flag checks &
\path{src/tensor_parallel_checks.py},
\path{reproducibility/tensor_parallel_sharding.md} \\
Deployment & Quantization, export, ONNX, and package-gate safety checks &
\path{src/quant_export_checks.py},
\path{reproducibility/quant_export_safety.md} \\
Developer UX & Human diagnostics, JSON, JUnit, SARIF, and GitHub annotation
reporters &
\path{src/reporters.py}, \path{tests/test_reporters.py} \\
Developer UX & Pre-commit hook, pytest plugin, GitHub Action path, and
baseline/suppression support &
\path{src/precommit.py}, \path{src/pytest_tensorguard.py} \\
Developer UX & Upstream-style \texttt{register\_forward\_pre\_hook} and
\texttt{@verifiable} decorator reference implementation &
\path{src/upstream_hook.py},
\path{reproducibility/upstream_hook_demo.md} \\
Documentation & Generated soundness and fragment specs plus concise onboarding &
\path{SOUNDNESS_CONTRACT.md}, \path{VERIFIABLE_FRAGMENT.md},
\path{GETTING_STARTED.md}, \path{LIMITATIONS.md} \\
Formal methods & Lean rules for core operators, reduced products, CEGAR,
fragment modes, and SMT encoding faithfulness &
\path{lean/TensorGuard/}, \path{tests/test_lean_whole_pipeline_audit.py} \\
Benchmarks & Frozen real benchmark corpus with datasheet and citation metadata &
\path{real_benchmarks/manifest.json}, \path{real_benchmarks/DATASHEET.md} \\
Benchmarks & Extended generated corpus, held-out blind split, false-positive
stress set, natural-distribution study, and mutation corpus &
\path{corpus_extended/}, \path{reproducibility/blind_split_eval.md} \\
Benchmarks & Public GitHub-mined bug snapshot and offline issue-miner proposal
flow &
\path{experiments_v5/github_bug_mining/}, \path{corpus_extended/issue_miner.py} \\
Evidence & Numeric-claim audit, dashboard, artifact capsule, benchmark delta
gate, and paper-evidence target &
\path{reproducibility/audit_numeric_claims.py},
\path{docs/dashboard/index.html}, \path{scripts/benchmark_delta.py} \\
\bottomrule
\end{longtable}

\section{Selected numerical claims}

\begin{center}
\small
\begin{tabular}{p{0.52\linewidth} p{0.32\linewidth}}
\toprule
\textbf{Claim} & \textbf{Backing artifact} \\
\midrule
117 operator confidence rows: 75 complete, 36 sound, 6 heuristic &
\path{operator_confidence_table.json} \\
16 frozen real benchmarks: 8 clean, 8 buggy &
\path{real_benchmarks/manifest.json} \\
2,704 GitHub-mined public bugs: 2,394 shape, 310 device &
\path{experiments_v5/github_bug_mining/mined_bugs_manifest.json} \\
TP=8, FP=0, TN=8, FN=0 on frozen real benchmarks &
\path{evaluation/confusion_matrices.json} \\
0 false positives on 80 sound-mode clean models &
\path{evaluation/sound_mode_fp.json} \\
0 false positives on 200 random clean fuzzed models &
\path{evaluation/diff_fuzz.json} \\
281/281 genuine injected faults caught &
\path{evaluation/neg_fuzz.json} \\
6/8 latent bugs caught vs. 0/8 for the strongest dynamic baseline &
\path{evaluation/hard_recall.json} \\
50 minimal bug reproducers frozen across nine mechanisms &
\path{evaluation/triage_regressions.json} \\
94.86\% frequency-weighted operator coverage on 13 torchvision models &
\path{evaluation/operator_frequency.json} \\
21/21 torchvision models traced and lowered by fx &
\path{evaluation/fx_trace_success.json} \\
371 public Lean theorems discovered and axiom-audited &
\path{tests/test_lean_whole_pipeline_audit.py} \\
12 verified, one regime-qualified, one environment-qualified numeric claims &
\path{reproducibility/numeric_claims_audit.json} \\
\bottomrule
\end{tabular}
\end{center}

\section{Manual bibliography}

\begin{thebibliography}{9}
\bibitem{pytorch}
PyTorch Contributors.
\newblock PyTorch: An imperative style, high-performance deep learning library.

\bibitem{z3}
Leonardo de Moura and Nikolaj Bjorner.
\newblock Z3: An efficient SMT solver.

\bibitem{cvc5}
The cvc5 Developers.
\newblock cvc5: A versatile and industrial-strength SMT solver.

\bibitem{pytea}
PyTea Contributors.
\newblock PyTea: Static analysis for tensor shape errors in deep learning
programs.

\bibitem{lean}
Lean Prover Community.
\newblock The Lean 4 theorem prover and programming language.

\bibitem{cegar}
Edmund Clarke, Orna Grumberg, Somesh Jha, Yuan Lu, and Helmut Veith.
\newblock Counterexample-guided abstraction refinement.

\bibitem{mcnemar}
Quinn McNemar.
\newblock Note on the sampling error of the difference between correlated
proportions or percentages.
\end{thebibliography}

\end{document}
